\documentclass[../../main.tex]{subfiles}
\begin{document}

$a$と$b$を互いに素，
すなわち1以外の公約数を持たない正の整数とし，
さらに$a$は奇数とする．
正の整数$n$に対して整数$a_n$，$b_n$を
${(a+b\sqrt{2})}^n=a_n+b_n\sqrt{2}$をみたすように定めるとき，
次の(1)，(2)を示せ．
ただし$\sqrt{2}$が無理数であることは証明なしに用いてよい．
\begin{enumerate}
  \item $a_2$は奇数であり，$a_2$と$b_2$は互いに素である．
  \item すべての$n$に対して，$a_n$は奇数であり，$a_n$と$b_n$は互いに素である．
\end{enumerate}

\end{document}
